% Generated by experiments/scripts/export_latex.py -- do not edit by hand.
% Source: experiments/data/2-processed-tables-and-evaluations/revised-usd-12m/holdout-w0/eval_summary.json
\begin{table}[htbp]
\centering
\caption[스펜더 홀드아웃 W0에서 짝지은 부트스트랩으로 구한 차이.]{스펜더 홀드아웃 W0에서 짝지은 부트스트랩으로 구한 차이예요(계약 $n=7{,}914$개, $a$와 $b$가 함께 쓰는 다시 뽑은 표본 400개, 시드 42). 첫 번째 묶음에는 규칙 A의 주된 부분에 들어가는 두 대비가 있어요. 이 두 대비는 이 연구가 자료를 하나도 뽑기 전인 2026년 10월 9일에 커밋한 분석 계획에서 미리 정했고, 여기서는 고친 점수에 적용했어요. 하나는 새 허락 쌍에서 C-PR과 C-PR ($\lambda=1$)을 비교하고, 다른 하나는 새 송금자에서 C-PR과 C-PR ($\lambda=0$)을 비교해요. 두 번째 묶음은 라벨을 본 뒤에 계산했어요. 여기에는 다음이 들어 있어요. 허락 층만 걷는 C-PR ($\lambda=1$)이 $t_1$에서의 자기 원시 차수보다 더 얻은 값(증가분), S-PR에 대한 같은 비교들, 그리고 엔도스랭크와 AWP가 각자의 원시 차수보다 더 얻은 값이에요. C-PR의 다른 대비는 표~\ref{tab:cpr-contrasts}에 있어요. 규칙 A의 주된 부분은 C-PR이 새 허락 쌍에서 C-PR ($\lambda=1$)보다 더 나쁘지 않고, 새 송금자에서 C-PR ($\lambda=0$)보다 더 나아야 한다는 것이에요. “더 나쁘지 않음”은 성립하지만 “더 나음”은 성립하지 않아요. 그래서 이 규칙은 충족되지 않아요. 엔도스랭크와 AWP 사이의 차이는 표~\ref{tab:endorserank-awp-holdout}에 있어요.}
\label{tab:holdout-diff}
\fitwidth{%
\begin{tabular}{llccccc}
\toprule
대비 ($a$ 빼기 $b$) & 라벨 & $\tau(a)$ & $\tau(b)$ & $\Delta\tau$ & $\Delta\tau$의 95\% 신뢰구간 & $P(\Delta\tau>0)$ \\
\midrule
\multicolumn{7}{l}{\emph{규칙 A, 주된 부분, 계획에서 미리 정함}} \\
C-PR 빼기 C-PR ($\lambda=1$) & 새 허락 쌍 & 0.329 & 0.298 & +0.031 & [0.021, 0.042] & 1.000 \\
C-PR 빼기 C-PR ($\lambda=0$) & 새 송금자 & 0.336 & 0.366 & -0.030 & [-0.040, -0.020] & 0.000 \\
\midrule
\multicolumn{7}{l}{\emph{라벨을 본 뒤에 계산함}} \\
C-PR ($\lambda=1$) 빼기 $t_1$ 받은 허락 수 & 새 허락 쌍 & 0.298 & 0.392 & -0.094 & [-0.104, -0.083] & 0.000 \\
S-PR 빼기 C-PR ($\lambda=1$) & 새 허락 쌍 & 0.190 & 0.298 & -0.108 & [-0.128, -0.088] & 0.000 \\
S-PR 빼기 C-PR ($\lambda=0$) & 새 송금자 & 0.377 & 0.366 & +0.011 & [0.004, 0.018] & 1.000 \\
EndorseRank 빼기 $t_1$ 받은 허락 수 & 새 허락 쌍 & 0.303 & 0.392 & -0.090 & [-0.100, -0.079] & 0.000 \\
AWP 빼기 $t_1$ 입차수 & 새 송금자 & 0.409 & 0.447 & -0.038 & [-0.043, -0.032] & 0.000 \\
\bottomrule
\end{tabular}
}
\end{table}
